We prove Proposition~\ref{prop:full-support-nonmonotonicity}. Throughout this
appendix,
\[
    k_t(x,y)=\exp\left\{-\frac{(x-y)^2}{2t}\right\};
\]
the omitted Gaussian normalizing constant is independent of the mixing
distribution. For a weighted empirical measure
\[
    P=\sum_{i=1}^q w_i\delta_{x_i},
    \qquad w_i>0,\quad \sum_iw_i=1,
\]
write \(q_{t,G}(x_i)=\int k_t(x_i,y)\,dG(y)\) and
\begin{equation}
\label{eq:nonmon-dual}
    D_{t,P,G}(y)
    =\sum_iw_i\frac{k_t(x_i,y)}{q_{t,G}(x_i)}.
\end{equation}
Lindsay's directional-derivative criterion
\cite{lindsay1983geometryA} says that \(G\) is an NPMLE exactly when
\begin{equation}
\label{eq:nonmon-kkt}
    D_{t,P,G}(y)\le1\quad(y\in\mathbb R).
\end{equation}
The identity
\begin{equation}
\label{eq:nonmon-dual-average}
    \int D_{t,P,G}(y)\,dG(y)=1.
\end{equation}
then makes equality automatic \(G\)-almost everywhere. Strict concavity of
the log likelihood in the fitted vector implies that all maximizers have the
same fitted vector. Thus, if the equality set in
\eqref{eq:nonmon-kkt} is finite and its kernel columns are linearly
independent, the NPMLE is unique.

We also use the standard stability of nondegenerate zeros: if \(F_N\to F\) in
\(C^1\) near \(z_0\), \(F(z_0)=0\), and \(DF(z_0)\) is invertible, then
\(F_N\) has a zero converging to \(z_0\).

The first lemma produces the high-variance dual certificate. Its first
\(r\) data locations remain fixed, while the last location and its contact
move to \(+\infty\).
\begin{lemma}
\label{lem:nonmon-remote-certificate}
Fix \(r\ge1\) and \(t>0\). There are fixed points
\[
    a_1<\cdots<a_r=0
\]
such that, for every sufficiently large \(B\), there are positive
coefficients \(c_1(B),\ldots,c_{r+1}(B)\) for which
\begin{equation}
\label{eq:nonmon-remote-certificate}
    C_B(y)
    =
    \sum_{i=1}^r c_i(B)k_t(a_i,y)
    +c_{r+1}(B)k_t(B,y)
    \le1.
\end{equation}
Equality holds at exactly \(r+1\) nondegenerate maxima
\(\theta_1(B),\ldots,\theta_{r+1}(B)\). Moreover,
\[
    c_i(B)\to c_i^0>0,\quad
    \theta_i(B)\to\alpha_i\quad(i\le r),
    \qquad
    c_{r+1}(B)\to1,\quad
    \theta_{r+1}(B)-B\to0,
\]
and the square contact matrix
\begin{equation}
\label{eq:nonmon-contact-matrix}
    K_B
    =
    \bigl(k_t(x_i,\theta_j(B))\bigr)_{i,j=1}^{r+1},
    \qquad
    (x_1,\ldots,x_{r+1})=(a_1,\ldots,a_r,B),
\end{equation}
is invertible.
\end{lemma}

\begin{proof}
We first build the \(r\) fixed peaks. Choose a large separation \(L\) and
put \(a_i=(i-r)L\). Let the coefficients
\(c=(c_1,\ldots,c_r)\) range over a fixed neighborhood of
\(\boldsymbol1\). For fixed \(R>0\), local coordinates \(y=a_i+s\) give,
uniformly over \(i,c\) and \(|s|\le R\),
\[
    \sum_{j=1}^r c_jk_t(a_j,a_i+s)
    =
    c_i e^{-s^2/(2t)}+E_{i,L,c}(s),
\]
where
\begin{equation}
\label{eq:nonmon-separated-error}
    \|E_{i,L,c}\|_{C^2[-R,R]}
    \le
    C_{r,R,t}(1+L^2)
    \exp\left\{-\frac{(L-R)^2}{2t}\right\}
    \longrightarrow0.
\end{equation}
Choose \(0<\delta<\min\{R,\sqrt t\}\). The derivative of the limiting
single peak is bounded away from zero on
\(\{\delta\le |s|\le R\}\), while its second derivative is negative on
\([-\delta,\delta]\). Hence each interval
\([a_i-R,a_i+R]\) contains exactly one critical point
\[
    \alpha_i(c,L)=a_i+s_i(c,L),
    \qquad s_i(c,L)\to0,
\]
which is a nondegenerate maximum.

Let \(M_i(c,L)\) be the value at this maximum. Since the spatial derivative
vanishes there,
\[
    \frac{\partial M_i}{\partial c_j}(c,L)
    =
    k_t(a_j,\alpha_i(c,L)).
\]
Consequently \(M(c,L)-\boldsymbol1\to c-\boldsymbol1\) in \(C^1\) near
\(\boldsymbol1\). Stability of nondegenerate zeros gives coefficients
\(c(L)\to\boldsymbol1\) for which all \(r\) maximum heights equal one.
Choose \(R\) so large that \(2r e^{-R^2/(2t)}<1\), and then choose
\(L>2R\) large enough that the coefficients \(c(L)\) and critical points
\(\alpha_i(c(L),L)\) constructed above exist. Outside the union of
the \(r\) intervals, every center is at distance at least \(R\), so the
Gaussian sum is strictly below one. Inside each interval its unique critical
point is its maximum, of height one. We have therefore constructed
\[
    H_0(y)=\sum_{i=1}^r c_i^0k_t(a_i,y)\le1
\]
with exactly \(r\) nondegenerate contacts \(\alpha_i\). The matrix
\[
    A=\bigl(k_t(a_i,\alpha_j)\bigr)_{i,j=1}^r
\]
is invertible because it converges to the identity as \(L\to\infty\).

We now add the remote peak. For coefficients near
\((c_1^0,\ldots,c_r^0,1)\), the old maxima persist near the
\(\alpha_i\), and a new nondegenerate maximum appears near \(B\). The map
from the coefficients to these \(r+1\) maximum heights converges in \(C^1\),
as \(B\to\infty\), to the decoupled height map. Its derivative at the
limiting coefficient vector is
\[
    \begin{pmatrix}
        A^{\mathsf T}&0\\
        0&1
    \end{pmatrix}.
\]
Stability of nondegenerate zeros gives positive coefficients for which all
\(r+1\) heights are one, with the asserted limits.

It remains to check the global inequality. On small fixed neighborhoods of
the old contacts, the perturbed old cluster has a unique strictly concave
maximum of height one, and the remote term is negligible. On the complement
of these neighborhoods in a large fixed interval, \(H_0\) has a fixed gap
below one. An identical local argument controls the new maximum in a fixed
neighborhood of \(B\). Finally, enlarge the old and remote regions until the
entire old cluster is below \(1/4\) outside the old region and the remote
term is below \(1/4\) outside the region around \(B\). This proves
\eqref{eq:nonmon-remote-certificate} globally and identifies every contact.
The matrix \(K_B\) converges blockwise to
\(\operatorname{diag}(A,1)\), so it is invertible for large \(B\).
\end{proof}

An equal-height certificate can be inverted into empirical weights. This
lets us prescribe positive masses at all high-variance contacts before
choosing the sample multiplicities.
\begin{lemma}
\label{lem:nonmon-inverse-kkt}
Suppose
\[
    C(y)=\sum_{i=1}^q c_i k_t(x_i,y)\le1
\]
has equality exactly at \(\theta_1,\ldots,\theta_q\), where every \(c_i>0\),
and suppose \(K=(k_t(x_i,\theta_j))_{i,j=1}^q\) is invertible. For
\(p\in\Delta_q^\circ\), define
\[
    M=\operatorname{diag}(c_1,\ldots,c_q)K,
    \qquad
    w=Mp,
    \qquad
    G=\sum_{j=1}^q p_j\delta_{\theta_j}.
\]
Then \(w\in\Delta_q^\circ\), and \(G\) is the unique NPMLE for the weighted
empirical distribution \(\sum_iw_i\delta_{x_i}\).
\end{lemma}

\begin{proof}
Every column of \(M\) sums to one because
\[
    \sum_iM_{ij}
    =\sum_ic_i k_t(x_i,\theta_j)
    =C(\theta_j)=1.
\]
Thus \(w\) is a strictly positive probability vector. Moreover,
\[
    q_{t,G}(x_i)=(Kp)_i,
    \qquad
    \frac{w_i}{q_{t,G}(x_i)}=c_i.
\]
The dual certificate of \(G\) is exactly \(C\), so
\eqref{eq:nonmon-kkt} proves optimality. Any other maximizer is supported on
the same contact set and has a mass vector \(p'\) satisfying \(Kp'=Kp\).
Invertibility gives \(p'=p\).
\end{proof}

We next show that a suitable two-cluster empirical distribution has exactly
two atoms at the smaller variance. The exponent condition below says that
the smaller cluster has much more mass than the Gaussian overlap between the
two clusters.
\begin{lemma}
\label{lem:nonmon-two-cluster}
Fix \(a_1<\cdots<a_r=0\). For \(B\to\infty\), let
\[
    P_B=\sum_{i=1}^r w_{i,B}\delta_{a_i}+w_{\star,B}\delta_B
\]
be a probability measure with all weights positive, and put
\(\varepsilon_B=w_{r,B}\). Suppose
\begin{align}
    \varepsilon_B&\to0,\label{eq:nonmon-eps-zero}\\
    \log\varepsilon_B
      &=-\alpha B^2+o(B^2)
        \quad\text{for some }0<\alpha<\frac12,\label{eq:nonmon-eps-rate}\\
    \frac{w_{i,B}}{\varepsilon_B k_1(a_i,0)}
      &\to0\qquad(i<r).\label{eq:nonmon-old-negligible}
\end{align}
Then, for all sufficiently large \(B\), the variance-one NPMLE is unique and
has exactly two atoms.
\end{lemma}

\begin{proof}
Condition \eqref{eq:nonmon-old-negligible} gives
\(\sum_{i<r}w_{i,B}=o(\varepsilon_B)\), and hence
\begin{equation}
\label{eq:nonmon-remote-weight}
    w_{\star,B}=1-\varepsilon_B\{1+o(1)\}.
\end{equation}
We seek an NPMLE of the form
\[
    G_{\rho,u,z}
    =
    \varepsilon_B\rho\,\delta_u
    +(1-\varepsilon_B\rho)\delta_{B+z},
\]
where \((\rho,u,z)\) is near \((1,0,0)\). Let
\begin{align*}
    q_i(\rho,u,z)
    &=
    \varepsilon_B\rho k_1(a_i,u)
    +(1-\varepsilon_B\rho)k_1(a_i,B+z),\\
    q_\star(\rho,u,z)
    &=
    \varepsilon_B\rho k_1(B,u)
    +(1-\varepsilon_B\rho)k_1(B,B+z),
\end{align*}
and denote the corresponding dual function by \(D_{\rho,u,z}\).

For every fixed \(C>0\) and integer \(j\ge0\),
\eqref{eq:nonmon-eps-rate} implies
\begin{equation}
\label{eq:nonmon-cross-cluster}
    B^j\frac{\exp\{-(B-C)^2/2\}}{\varepsilon_B}\longrightarrow0.
\end{equation}
Indeed, the logarithm of the left side is
\(-(\frac12-\alpha)B^2+O(B)+o(B^2)\). Since every derivative of a Gaussian
kernel is the kernel times a polynomial,
\eqref{eq:nonmon-cross-cluster} controls all parameter and spatial
derivatives of the cross-cluster terms below.

Set \(\lambda_i=w_{i,B}/q_i\) and
\(\lambda_\star=w_{\star,B}/q_\star\). Uniformly for
\((\rho,u,z)\) near \((1,0,0)\),
\begin{align}
    \lambda_i&\to0
      &&\text{in \(C^1\), }i<r,\label{eq:nonmon-lambda-small}\\
    \lambda_r&\to\frac{1}{\rho k_1(0,u)}
      &&\text{in \(C^1\)},\label{eq:nonmon-lambda-zero}\\
    \lambda_\star&\to\frac{1}{k_1(0,z)}
      &&\text{in \(C^1\)}.\label{eq:nonmon-lambda-star}
\end{align}
For example, \(q_i\ge\varepsilon_B\rho k_1(a_i,u)\), so
\[
    0\le\lambda_i
    \le
    C\frac{w_{i,B}}{\varepsilon_Bk_1(a_i,0)}
    \longrightarrow0.
\]
Differentiation introduces either a bounded local derivative or a
cross-cluster derivative divided by \(\varepsilon_B\), which is controlled
by \eqref{eq:nonmon-cross-cluster}. Differentiating the identity
\(w_{r,B}=\varepsilon_B\) in the same way proves
\eqref{eq:nonmon-lambda-zero}. Finally,
\eqref{eq:nonmon-lambda-star} follows from
\eqref{eq:nonmon-remote-weight} and
\[
    q_\star
    =
    (1-\varepsilon_B\rho)k_1(0,z)
    +\varepsilon_B\rho k_1(B,u).
\]

Consider the three equations
\[
    F_B(\rho,u,z)
    =
    \begin{pmatrix}
        D_{\rho,u,z}(u)-1\\
        D'_{\rho,u,z}(u)\\
        D'_{\rho,u,z}(B+z)
    \end{pmatrix}
    =0.
\]
Equations
\eqref{eq:nonmon-cross-cluster}--\eqref{eq:nonmon-lambda-star} imply
convergence in \(C^1\) near \((1,0,0)\) to
\[
    F_\infty(\rho,u,z)
    =
    \begin{pmatrix}
        \rho^{-1}-1\\
        -u/\rho\\
        -z
    \end{pmatrix}.
\]
Since \(F_\infty(1,0,0)=0\) and
\(DF_\infty(1,0,0)=-I_3\), there is a solution
\[
    \rho_B\to1,\qquad u_B\to0,\qquad z_B\to0.
\]
Put
\[
    \pi_B=\varepsilon_B\rho_B,
    \qquad
    v_B=B+z_B,
    \qquad
    G_B=\pi_B\delta_{u_B}+(1-\pi_B)\delta_{v_B}.
\]
For all large \(B\), \(0<\pi_B<1\), and construction gives
\[
    D_B(u_B)=1,\qquad D_B'(u_B)=D_B'(v_B)=0.
\]
The second contact height is automatic: by
\eqref{eq:nonmon-dual-average},
\[
    1=\pi_BD_B(u_B)+(1-\pi_B)D_B(v_B),
\]
so \(D_B(v_B)=1\).

It remains to exclude other contacts. At the solution,
\[
    \lambda_r\to1,\qquad
    \sum_{i<r}\lambda_i\to0,\qquad
    \lambda_\star\to1.
\]
Consequently,
\[
    D_B(y)\to k_1(0,y)
    \quad\text{in \(C^2\) on fixed compact sets},
\]
and
\[
    D_B(B+s)\to k_1(0,s)
    \quad\text{in \(C^2\) on fixed compact \(s\)-sets}.
\]
Choose \(0<\delta<1\). The limiting Gaussian is strictly concave on
\([-\delta,\delta]\), so \(u_B\) is the unique contact there. Uniform
convergence gives a strict gap below one on
\([-R,R]\setminus(-\delta,\delta)\). The translated argument shows that
\(v_B\) is the only contact in \([B-R,B+R]\). Finally choose \(R\) so large
that \(3e^{-R^2/2}<1\). Outside these intervals,
\[
    D_B(y)
    \le
    (1+o(1))k_1(0,y)
    +(1+o(1))k_1(B,y)+o(1)
    \le2e^{-R^2/2}+o(1)<1.
\]
Thus \(G_B\) satisfies \eqref{eq:nonmon-kkt}, with equality exactly at its
two atoms.

The two kernel columns are linearly independent because the rows
corresponding to \(0\) and \(B\) have determinant
\[
\begin{aligned}
    &\det
    \begin{pmatrix}
        k_1(0,u_B)&k_1(0,v_B)\\
        k_1(B,u_B)&k_1(B,v_B)
    \end{pmatrix}\\
    &\qquad
    =
    k_1(0,v_B)k_1(B,u_B)
    \left\{e^{B(v_B-u_B)}-1\right\}>0.
\end{aligned}
\]
The NPMLE is therefore unique.
\end{proof}

We now choose the high-variance fitted masses so that the induced empirical
weights satisfy the two-cluster lemma at variance one.
\begin{proof}[Proof of Proposition~\ref{prop:full-support-nonmonotonicity}]
Rescaling the observations and mixing locations by \(S^{-1/2}\) reduces the
claim to \(S=1\) and \(T=\tau>1\). Put \(r=m-1\), and apply
Lemma~\ref{lem:nonmon-remote-certificate} at variance \(\tau\). Thus
\[
    x_1=a_1<\cdots<x_r=a_r=0<x_m=B
\]
support a certificate
\[
    C_B(y)=\sum_{i=1}^mc_i(B)k_\tau(x_i,y)\le1
\]
with exactly \(m\) contacts, the last of which satisfies
\(\theta_m(B)=B+o(1)\). Let \(K_B\) be the contact matrix
\eqref{eq:nonmon-contact-matrix}, and set
\[
    M_B=\operatorname{diag}\{c_1(B),\ldots,c_m(B)\}K_B.
\]

Choose the high-variance NPMLE masses
\[
    \eta_B=e^{-B^2/\tau},
    \qquad
    p_j^{(B)}=\frac{\eta_B}{r}\quad(j\le r),
    \qquad
    p_m^{(B)}=1-\eta_B,
\]
and define the empirical weights \(w^{(B)}=M_Bp^{(B)}\).
Lemma~\ref{lem:nonmon-inverse-kkt} gives a unique \(m\)-atom NPMLE at
variance \(\tau\).

For \(i\le r\), write
\[
    q_i^{(B)}
    =
    \sum_{j=1}^mp_j^{(B)}k_\tau(a_i,\theta_j(B)).
\]
Since the first \(r\) masses sum to \(\eta_B\),
\[
    q_i^{(B)}
    =
    (1-\eta_B)k_\tau(a_i,\theta_m(B))+O(\eta_B).
\]
Because \(a_i\) is fixed and \(\theta_m(B)=B+o(1)\),
\[
    \frac{\eta_B}{k_\tau(a_i,\theta_m(B))}
    =
    \exp\left\{-\frac{B^2}{2\tau}+O(B)+o(B)\right\}
    \longrightarrow0.
\]
It follows that
\begin{equation}
\label{eq:nonmon-q-dominance}
    q_i^{(B)}
    =
    \{1+o(1)\}k_\tau(a_i,\theta_m(B)).
\end{equation}
Put \(\varepsilon_B=w_r^{(B)}=c_r(B)q_r^{(B)}\). Since
\(c_r(B)\to c_r^0>0\),
\[
    \log\varepsilon_B
    =
    -\frac{B^2}{2\tau}+o(B^2).
\]
Thus \eqref{eq:nonmon-eps-rate} holds with
\(\alpha=1/(2\tau)<1/2\). For \(i<r\),
\begin{align*}
    \frac{w_i^{(B)}}{\varepsilon_Bk_1(a_i,0)}
    &=
    \{1+o(1)\}\frac{c_i^0}{c_r^0}
    \frac{k_\tau(a_i,\theta_m(B))}
         {k_\tau(0,\theta_m(B))k_1(a_i,0)}\\
    &=
    \{1+o(1)\}\frac{c_i^0}{c_r^0}
    \exp\left\{
        \frac{a_i\theta_m(B)}{\tau}
        +\frac{a_i^2}{2}\left(1-\frac1\tau\right)
    \right\}
    \longrightarrow0,
\end{align*}
because \(a_i<0\) and \(\theta_m(B)\to\infty\).
Lemma~\ref{lem:nonmon-two-cluster} now gives a unique two-atom NPMLE at
variance one for all sufficiently large \(B\).

It remains to replace the weights by integer multiplicities. The matrix
\(M_B\) is invertible and satisfies
\(\boldsymbol1^{\mathsf T}M_B=\boldsymbol1^{\mathsf T}\). Hence
\(M_B\Delta_m^\circ\) is relatively open in the affine probability
hyperplane and contains \(w^{(B)}\). Along a sequence \(B_N\to\infty\),
choose a rational vector
\[
    \widetilde w^{(N)}
    \in M_{B_N}\Delta_m^\circ\cap\mathbb Q^m
\]
so close to \(w^{(B_N)}\) that
\[
    \max_i
    \left|
       \frac{\widetilde w_i^{(N)}}{w_i^{(B_N)}}-1
    \right|\longrightarrow0.
\]
Then
\(\widetilde p^{(N)}=M_{B_N}^{-1}\widetilde w^{(N)}\) lies in
\(\Delta_m^\circ\), so the same certificate gives a unique \(m\)-atom NPMLE
at variance \(\tau\). The relative approximation preserves
\eqref{eq:nonmon-eps-rate} and \eqref{eq:nonmon-old-negligible}, so the
variance-one NPMLE still has two atoms. Expressing the coordinates of one
such rational vector over a common denominator gives an ordinary sample with
exactly \(m\) distinct values and positive integer multiplicities. Rescaling
by \(\sqrt S\) completes the construction.
\end{proof}
